\section{Model interfaces and qualification}\label{model-interfaces-and-starter-architectures}

Launch admits one \textit{A} model per entry. The model predicts AI share,
production history, author identity, and AI spans in one invocation. All output heads
must consume a common document encoding produced by a shared learned encoder;
each may also have task-specific layers. This requirement enforces one shared
encoding interface; whether the heads use common predictive features requires
measurement. Admission checks this structure in the
submitted computation graph, together with one manifest and input/output
contract. Packaging two independent predictors behind one API does not qualify.
This checks inference structure, not whether the components were trained jointly.
Known AI models have predefined author profiles alongside human
writers. Detecting AI involvement is the primary purpose of the origin output
and the only reward quantity; the other outputs are qualification gates. The
transformation mechanism \textit{B} is deferred to Section~\ref{later-b};
\textit{C}'s proposed contributions to the shared grammar that \textit{B} uses
appear in Section~\ref{grammar-contributions}.

The shared encoder is intended to let the two tasks improve each other.
Learning author- and generator-specific writing patterns can help identify
model involvement. Production-history supervision can help attribution
distinguish a writer's habits from patterns introduced by generation or
editing. Both objectives can therefore shape features useful to the other.
This expected mutual benefit requires measurement: sharing an encoder alone
does not establish it. Section~\ref{development-and-launch-criteria} requires
each new architecture to be tested in both directions against models trained
for each task alone.
The requirement deliberately gives detection, attribution, and later author-fit
feedback one reusable learned representation. Graph inspection cannot establish
useful feature sharing; disjoint slices could satisfy the structural rule.
Measured harm would prompt an explicit revision of the design before launch,
rather than an undocumented exception to admission.

An \textit{A} entrant submits the weights and all data needed for the declared
inference calculation. The artifact stays private until its cohort's scoring
starts and is published then (Section~\ref{a-artifact-execution}). Miner
submissions are licensed CC BY-SA 4.0, so anyone may obtain a published artifact
free of charge and execute it locally, and later \textit{B} development is
expected to train against published detectors. A published artifact stays
public and is never recalled. The coordinator deletes its private admission
copy once the cohort closes and publication is recorded, and keeps the current
and previous reference artifacts hosted for evaluation while they are
references (Section~\ref{cohort-receipt-retention}). Hosted execution through
the paid API is charged separately.
Miners may keep their training recipes and permitted training corpora private.

Project code and documentation are licensed MIT OR Apache-2.0, except files and
datasets that carry their own licence or notice. Reference weights are licensed
CC BY-SA 4.0, like miner submissions, and Slop Ninja publishes its own
references, \(R_0\) and any later reference it trains
(Section~\ref{reference-model}), on the same terms. A public reference package
carries no evaluation text: its replay cases are synthetic rather than
development documents, and the package passes a scan for corpus text before
release. \(R_0\) is published once its package has been re-exported in this
way. Miners can also start from a public training recipe~\cite{trainingguide}:
it fine-tunes ModernBERT-base on about 4,200 Train rows from five
redistributable dataset snapshots and packs and checks an admissible artifact
with the miner client. It reproduces \(R_0\)'s pipeline and artifact format,
not its V6 starting point, data, or weights, and is not a strong detector. An
entrant artifact that becomes the reference was already published with its
cohort.
Slop Ninja executes submitted artifacts for the pilot benchmark and published,
qualified artifacts for customer service; an \textit{A} miner needs no private
serving account or delegated inference credential.

\subsection{AI detection through origin and attribution}\label{author-and-origin-detection}

The origin output is a calibrated probability distribution over three recorded
production histories: human-only (\(p_{\mathrm{human}}\)), model-only
(\(p_{\mathrm{model}}\)), and mixed human/model work (\(p_{\mathrm{mixed}}\)). The
artifact combines model-only and mixed probabilities to answer the primary
AI-involvement question. It reports a continuous probability, displayed as a
percentage:
\[
 p_{\mathrm{AI}} = p_{\mathrm{model}} + p_{\mathrm{mixed}},
 \qquad
 \texttt{ai\_confidence\_percent} = 100\,p_{\mathrm{AI}}.
\]
Binary training labels are zero for documented human-only work and one for
model-only or mixed work. The model learns a probability from those labels;
it does not convert a hard decision into a percentage. Documented three-way
histories additionally supervise the origin components. The three probabilities
sum to one, and their calibration method and workload profile are frozen with
the artifact. Every completed inference returns numeric share and confidence
scores. An unavailable prediction returns an explicit failure status, never
zero scores. Fixed
thresholds remain useful for benchmark operating points but are not required
public outputs.

The second percentage is estimated AI share: the expected fraction of final
words in model-written or model-rewritten spans under the provenance convention
in Section~\ref{training-data-and-evaluation-separation}. Write that fraction
as \(f_{\mathrm{AI}}\) and its predicted expectation as \(\hat f_{\mathrm{AI}}\):
\[
 \texttt{ai\_share\_percent}=100\,\hat f_{\mathrm{AI}},
 \qquad 0\leq p_{\mathrm{model}}\leq\hat f_{\mathrm{AI}}\leq p_{\mathrm{AI}}\leq1.
\]
These bounds follow because human-only work has zero share and model-only
work has unit share under this convention. One implementation uses
\(\hat f_{\mathrm{AI}}=p_{\mathrm{model}}+p_{\mathrm{mixed}}\mu_{\mathrm{mixed}}\),
where \(\mu_{\mathrm{mixed}}\in[0,1]\) is the learned expected share conditional
on mixed history and the document. Binary labels supervise confidence; measured
fractional labels additionally supervise share. A result can therefore report
20\% estimated share with 99\% confidence that AI contributed. Confidence is
the probability of involvement, not certainty that the share estimate is exact;
uncertainty intervals for share require separate validation. Both percentages
must qualify on their declared workloads before paid service.

Attribution reports conditional results for both human and AI-model roles.
Each result names supported gallery matches, or returns \emph{Other human}
(\texttt{other\_human}) or \emph{Other AI} (\texttt{other\_ai}) when no named
profile for that role satisfies the frozen matching rule. The question is
which contributor fits \emph{if that role contributed}; a conditional match or
fallback does not itself assert that the role is present. Report its involvement
probability alongside it:
\[
 p_{\mathrm{human\ involved}}=p_{\mathrm{human}}+p_{\mathrm{mixed}},
 \qquad p_{\mathrm{AI\ involved}}=p_{\mathrm{AI}}.
\]
These probabilities may sum above one because mixed work includes both roles;
they are not shares of the text. No presence cutoff suppresses a role's
attribution. Both fallbacks can therefore appear for mixed work, including
when neither contributor has a named match. Multiple contributors need not
be forced into one identity.

The fallback categories are outcomes, not enrolled profiles or assertions
that an identity is absent from the gallery. A listed writer or model may be
unidentifiable from the available text. There is no separate unknown or
attribution-abstention label in a completed result; failure remains separate.
Supported named-model matches remain visible even at low AI-involvement
probability, without changing that probability or proving authorship.
The output contract fixes these meanings and their calibration before evaluation.

Origin and attribution answer related but different questions. An unseen generator
may receive high AI confidence and \emph{Other AI} attribution. Failure to match a
known model does not establish human origin. A known writer's style may survive
AI assistance, so a human-profile match cannot override the origin estimate.
When a model rewrites human-only text, the resulting work has mixed history.
When one model rewrites another model's draft, it remains model-only under the
initial label policy, with attribution potentially unresolved. The outputs
estimate the recorded workflow and its contributors; they do not prove authorship.

Independent reference works form the gallery. Its construction, pooling,
reference/query roles, and version are part of the inference contract. The
interface should permit enrollment of a new human author from independent
samples without fitting another permanent identity classifier. Profile IDs
identify output matches and gallery entries; they are excluded from query
features.

The main intended human-attribution workflow is personal: a customer supplies
their own writing samples to construct a private writer profile. A qualifying
match returns \emph{Provided writer} (\texttt{provided\_writer}), bound to that
customer's selected reference profile. This is a named match under the normal
matching rule; if no human profile qualifies, return \emph{Other human}. It can accompany
AI involvement and identifies resemblance to the supplied writing, without
authenticating the customer's identity or proving authorship. The profile is
scoped to the request or customer account and is never implicitly enrolled in
the public gallery. Operator-maintained human profiles may also use documented
public or licensed references under a published enrollment policy recording
provenance, the identity basis, and permitted uses. Public availability alone
does not settle those uses.

In this later workflow, trusted preprocessing represents each reference work by
\textit{A}'s encoder output and by stylometric features, function-word
frequencies and character n-grams, then pools them into one writer profile under
a qualified, versioned pooling contract. A normalized weighted mean is the
initial candidate; the contract fixes feature scaling, per-work weights,
missing-channel handling, and the compatible encoder, feature, and pooling
versions. The
matcher receives the profile instead of raw reference prose. A pooled summary
need not represent any single work, so qualification must test whether
averaging loses useful variation across topics or registers and whether sample
length or an outlier dominates the result. Raw samples are still read during
trusted preprocessing. Both samples and profiles receive customer-content
protection; a profile is not an anonymization method. The same fixed profile is
the intended target for later \textit{B} revision using \textit{A}'s continuous
author-fit feedback (Section~\ref{author-transformation-and-detector-evasion}).

No profile type yet qualifies to back \emph{Provided writer} or that feedback.
In an October 31, 2026 comparison on writers neither system had seen, unlearned
profiles from \(R_0\)'s encoder beat CFG writer vectors on verification,
stylometric features beat both, and adding CFG to the encoder plus stylometry
gained only \(+0.0015\) in AUC, so CFG writer vectors are not a profile
channel. Encoder plus stylometry was the strongest
profile, yet at five references 42\% of impostor trials still passed when 90\%
of genuine ones did, and no profile ranked rewrites toward a target writer much
above chance (Section~\ref{writer-profile-evidence}). The next qualification is
a learned matcher trained on fresh authors, with its acceptance thresholds
measured on withheld writers and on real style-transfer data. Until a profile
type passes it, neither \emph{Provided writer} nor \textit{B}'s author-fit
feedback is offered.

Open-set rejection is evaluated on withheld writers and
generators, separately from choosing among known profiles. Rejection means a
completed attribution decision that accepts no named match for the evaluated
writer or generator role, including \emph{Other human} and \emph{Other AI}.
Evaluate each documented contributor role, including both roles in mixed work,
regardless of its predicted probability. Credit successful rejection only when
the ground-truth identity is absent from the gallery and the corresponding
fallback accepts no named match for that role. A wrong named match is an
attribution error; fallback on a listed identity is not successful open-set
rejection. Missing role results and execution failures remain in the assigned
population and earn no rejection credit. Do not score conditional identities
for roles absent from the documented history; evaluate mistaken involvement
probabilities through origin calibration and false-positive measurements.
This role-level rejection rule covers one contributor or a role whose
contributors are all withheld. A role containing both known and withheld
contributors needs a frozen set-valued attribution metric before earning
rejection credit; retain such cases in coverage and report them separately.
Launch qualification covers English documents. Multilingual claims and
additional document-length classes need their own data and workload evidence.

\subsection{AI span localization}\label{span-outputs}

The goal is to locate AI-written text, not only to score whole documents.
Sentences are one view of that location. The output contract therefore adds a
token-level head on the shared encoder. It returns an AI probability for every
token, decodes those probabilities into AI spans with exact character offsets
into the canonical text and a confidence for each span, and derives per-sentence
AI share and AI-involvement confidence from the same token probabilities. The
document-level outputs keep their meanings, and the span outputs are made to
agree with them by construction. In the runtime's emission, a word's
probability is the mean probability of its tokens, and a share tilt shifts every
word's logit by one common bias so that the mean word probability equals the
emitted document share. A sentence's
share is the mean tilted probability of its words, so sentence shares aggregate
to the document share by final-word counts over complete, nonoverlapping
partitions. Document involvement confidence cannot be assembled that way,
because neighboring sentences are dependent: neither averaging sentence
confidences nor multiplying their complements yields a document probability.
It remains the document head's calibrated output, and sentence and span
confidences are capped at it.

The training labels already exist: retained exact model-produced spans for mixed
documents, all-AI labels for model-only documents, and all-human labels for
human-only documents (Section~\ref{training-data-and-evaluation-separation}).
Qualification measures token-level Brier score, character-span intersection over
union and F1 against the retained spans, span-boundary error, and sentence-level
calibration. Each cohort reports these measures as a qualification gate whose
thresholds are fixed after measurement; span outputs earn no reward, and
document involvement Brier remains the reward. Agreement with the document
outputs needs no gate, because the emission guarantees it; the gap between mean
word probability and document share before the tilt is reported as a
diagnostic, since the tilt costs word-level accuracy wherever the document
share is the less accurate estimate. The span gate includes the
false-positive rate of human sentences inside mixed documents, reported
separately from the pure-human document gate, because identical human words
score higher next to AI text than alone (Section~\ref{cohort-mix-evidence});
like every span measure it carries no reward weight.

The release package pairs \(R_0\), the first reference model, with such a head
(Section~\ref{span-evidence}). The head was
trained on \(R_0\) with the encoder and document heads frozen, so its document
outputs are bit-identical to \(R_0\)'s. Its decoder was chosen under a limit of
spans on at most 22 of 570 development human documents, which it meets with no
slack. On held-out spliced rewrites it reaches mean per-document character F1
0.410 and recall 0.402, counting a document without a decoded span as zero, and
precision 0.862 over the documents in which it decodes one.
It decodes no span in 42\% of those documents and mostly misses clause and
single-sentence edits, at character F1 0.051 and 0.167. Retrospectively, on the
consumed sealed Test, 2.5\% of human documents carry a decoded span, and text
written by a base model is almost never found. A joint fine-tune of encoder and
heads found more spans but scored human documents clearly worse before
calibration, so it was not adopted.

Slop Ninja has adopted a follow-up head, S2c, as its span head for the first
package that serves spans; until span serving lands, the shipped head above
stays in the release package. S2c keeps the encoder and document heads frozen
and feeds the token head through a trained copy of the encoder's last four
layers and final norm. Its document outputs are byte-identical to \(R_0\)'s at
every seed, so confidence, share, and reward scoring do not change. Its own
decoder puts spans on 21 of the 570 development human documents, within the
limit, though seeds 29 and 41 give 22 and 27. Compared without span calibration
maps, on the held-out spliced rewrites it reaches character F1 0.519 and span
F1 0.485, against 0.410 and 0.241 for the shipped head (0.245 with its maps);
retrospectively, on the consumed Test, its mixed-document character F1 is 0.535
against 0.478, but with calibration 3.4\% of human documents carry a span
against 2.5\% (Section~\ref{s2c-evidence}). The tower adds 18--20\% of the
encoder's CPU time and about 20 million parameters, which puts such a graph
outside both allow-listed graphs and above the 160 million parameter limit
(Section~\ref{a-artifact-execution}); admitting such a graph needs a versioned
manifest that widens both. No calibrated package with S2c has been built or
replayed yet. \(R_0\) remains testnet's first reference, and no span package
is named \(R_0\) (Section~\ref{reference-model}).

Until a cohort manifest enforces the span gate, the token-level head is
optional and its outputs are reported without affecting eligibility; from that
manifest on, every entrant must provide it. The allow-list names the graph with
the token-level head from launch, but the scorer does not run span heads yet,
so admission refuses that graph until span serving qualifies
(Section~\ref{a-artifact-execution}). Enforcing the gate does not by itself
change the reference; reference changes follow Section~\ref{reference-model}.

A further extension could identify words and constructions that materially
influence a score, starting from the token-level outputs. Test explanations by
removing or changing the cited text and measuring the model's response with
meaning and context controls. Require stability across innocuous formatting
changes. Explanations in terms of grammatical constructions belong to later
\textit{B} revision guidance, which uses the shared grammar. An
influence on a model score is evidence about its calculation, not proof that a
word or construction was produced by AI. Attention weights or a generated
explanation alone do not establish that influence. These explanations are a
separate future qualification, not an initial launch requirement.

\subsection{\textit{A} input contract: raw text}\label{input-contract}

The \textit{A} input is the canonical document text. A miner artifact is a
text model: an ONNX computation graph, a fixed tokenizer, the output heads, and
their calibration. The coordinator tokenizes each document with the
allow-listed tokenizer, which the artifact's own must match byte for byte, and
runs the graph on CPU in a Slop Ninja-controlled Cloudflare
Container under the cohort's declared length and resource limits; an input beyond the artifact's limit returns an explicit
unsupported-length result. Cohort batches stay within the declared profile.
While the span gate is inactive, an entrant with document heads only, like
\(R_0\), is admissible (Section~\ref{span-outputs}).
The reference \(R_0\) is a ModernBERT-base encoder~\cite{modernbert} of
about 150M parameters with a 2,048-token limit, fine-tuned from the earlier V6
detector on raw tokens, with origin, share, and attribution heads
(Section~\ref{existing-evidence-and-implementation-status}); the release
package adds the token-level span head. Its architecture,
with or without a token-level span head, defines the graphs that admission
allows (Section~\ref{a-artifact-execution}).

\textit{A} uses no features derived from a resolved context-free-grammar (CFG)
parse. A CFG-only detector failed its first real-data test: on an October 26,
2026 research Test of 102 documents
(splits as in Section~\ref{splits}),
CFG-only shared models had AI-involvement Brier 0.126--0.134 against 0.096 for a
Calibration-mapped Pangram comparator on the same documents, and flagged
67--87\% of human documents at the 0.5
cutoff~\cite{cfgviability}. Adding CFG coordinates to frozen text features gave
only a small, uncertain Development gain, 0.048--0.057 against 0.052--0.058 for
text alone~\cite{v6cfgfusion}. Those tests used pooled grammatical-leaf counts;
ordered and tree-aware CFG inputs remain untested. Parsing and ambiguity
resolution also add latency and cost, and numeric features carry information
about the source, so they provide no confidentiality guarantee.

The grammar's planned uses are \textit{B}'s edit representation, which types
edits and places a revision alongside its source (Section~\ref{later-b}), and
\textit{C}'s contributions to it (Section~\ref{grammar-contributions}); neither
depends on \textit{A}'s input contract. Any change to that contract, including
grammar features, would need a new paired comparison against raw-text models,
including its added cost, before admission
(Section~\ref{development-and-launch-criteria}).

\subsection{Artifact submission and execution policy}\label{a-artifact-execution}

An execution manifest fixes the whole calculation: weights, input/output
schema, feature meanings and order, normalization, availability masks,
serialization, calibration, reference pooling, numerical precision, rounding
and tolerances, admitted operators, dependencies, resource bounds, and
reference outputs. It also fixes the tokenizer and the canonical text profile.
Names, handles, and other identity-bearing material follow the profile's
declared policy. Explicit author IDs, miner UIDs,
assignment metadata, hidden labels, and source/candidate-role flags are
excluded from learned query inputs.

Models within a comparison receive the same accepted text. A change to the
input profile begins a new comparison and requires requalification;
it cannot silently alter an unfinished cohort. Weight, calibration, or other
inference-semantic changes create a new artifact version. Compatible
operational repairs may preserve those semantics under the cohort rules in
Section~\ref{cohorts}. In particular, replacing an execution host must preserve
the artifact and calculation rather than substitute a different detector.

The coordinator caches the cohort batch, executes qualified artifacts under
equal limits, and retains their predictions and failures. Centrally executing
submitted artifacts keeps hidden benchmark inputs out of miner request paths
and limits direct benchmark harvesting. Model isolation and operator access
controls remain necessary. Publishing each executable when its cohort's scoring
starts permits inspection and independent experiments.
Validator duties follow Section~\ref{participants-and-the-two-kinds-of-work};
the scoring trust assumption is stated in
Section~\ref{trust-assumptions-and-adversarial-behavior}.

The manifest's allow-list names two graphs: the reference's ModernBERT-base
architecture and the same architecture with the token-level span head, each
with matching operators and shapes and integer arithmetic on token identifiers
only in the embedding lookup. An artifact has at most 160 million
parameters; \(R_0\) has 149.6 million. Its tokenizer must match the
allow-listed tokenizer byte for byte. Its calibration, contract, and other
declared files must validate against fixed schemas for their roles, such as a
monotone piecewise-linear calibration map whose knot count the schema bounds,
and a file without a declared role is rejected. These rules bound the total
size at about 700 MB, against the 610.6 MB of \(R_0\)'s artifact object, whose
graph file is 601 MB, so no separate byte cap
applies. The scorer serves only the graph without the span head today, so
admission refuses the span graph, recognized from the graph file's length in
the artifact object's header (below), until span serving
qualifies~\cite{artifactspec}. Where a limit is needed it derives in this way from structure, such as
the allow-list, a schema, or a measured bound, and is versioned in the manifest
rather than hard-wired, because a fixed limit that later models outgrow can
strand a subnet. The allow-list widens only through a versioned manifest. The
cap bounds what an artifact can memorize, an estimated 40--72 MB of learned
capacity, and an embedded index or lookup table
of text fails admission (Section~\ref{lookup-exposure}). The scorer has no
network access, online learning, cross-job prompt cache, or undeclared mutable
state. Qualification exercises resource bounds and numerical
boundary cases on the actual admitted runtime. Linux and macOS tooling may
inspect and prepare artifacts, but outputs from another backend become
interchangeable only after numerical qualification.

Qualification uses independently labelled hidden material and includes
formatting variants, identity and topic leakage, targeted triggers, chosen-input
failures, exact or near-copy submissions, and batch invariance: each document's
outputs must be identical whether it is scored alone or among its batch. Canonicalization does not replace
these controls. A publicly reproducible calculation can still exploit a
benchmark shortcut or favor a collaborator. These checks inform qualification
and continuing evaluation without establishing universal robustness.

\paragraph{Artifact objects.}
An artifact object is a length-prefixed header of at most 1 MiB, listing each
file's path, role, SHA-256, and length, followed by the files in path order.
The graph, weights, tokenizer, calibration, and output-schema files form its
inference content; reference test vectors, owner metadata, and publisher
signatures do not. The inference-content digest covers the frozen contract and
each inference file's role, digest, and length, so renaming or re-signing an
artifact cannot create new inference content. Two inference files with the same
role, digest, and length are refused. The object is the miner's local file and
its published form, but it is uploaded and published as the header and each
distinct file separately, from which anyone can rebuild it. The miner hashes
each file when it packs the artifact, and R2 computes each object's SHA-256 as it
arrives to check the declared digest; no Slop Ninja code hashes a model byte
after packing: the graph's class is checked by comparing its structure, the graph
file without its weight values, byte for byte with the allow-listed class's
pinned reference bytes~\cite{artifactspec}.

\paragraph{Private submission and duplicate credit.}
A miner commits on chain only a digest. Through the chain's commitment call, it
publishes a domain-separated commitment binding the chain and subnet, the
cohort, its hotkey and registration generation, the inference-content digest,
and a fresh secret 256-bit nonce. The commitment binds the inference content,
not the object's bytes, so objects that differ only in excluded files open the
same commitment. Before the admission cutoff the miner opens a signed admission
session with its commitment opening and the header's length and SHA-256, then
uploads the header and each file straight to the coordinator's private R2
admission storage, one presigned upload per object~\cite{cfr2}. Each upload URL
binds the object's key, its declared length and SHA-256, and create-only
semantics, so R2 verifies the bytes as they arrive and stores nothing on a
mismatch; no upload passes through a Slop Ninja Worker. The header's length and
SHA-256 are declared in the signed session request, and each file's are listed
in the header. File URLs are issued only after the header has arrived and the
coordinator has recomputed the inference-content digest from it and checked it
against the committed one, so the commitment pins each inference file's digest
and length; the header's own digest and its entries for excluded files are
declared, not committed. Admission, publication, and scoring then trust R2's stored
checksums and sizes and never hash model bytes again. TLS and the bucket's
access controls protect an artifact until publication; no further encryption
layer is applied, because R2 can verify only the bytes it receives. The
coordinator checks that the opening opens a finalized commitment for this
cohort on this chain, made by the hotkey and registration generation the
opening names, registered no later than the commitment block, that R2's stored
digests match the header, and that the artifact qualifies. That commitment's
finalized block and extrinsic index fix the entrant's admission order. A
commitment without a matching upload reserves nothing. Artifacts stay private
evaluation material until the cohort's scoring starts, after its roster is
frozen and its batch sealed, so no other participant can download, run, or copy
a model before it is entered; a detected copy entered in a later cohort earns
nothing (Section~\ref{public-copies}). When scoring starts the coordinator publishes
every admitted artifact, header and files each named by its SHA-256 and copied
inside R2 or by a stream that R2 verifies, and a published artifact is never
recalled. The commitment openings are published when the cohort closes. Other
validators need none of these objects to commit: they need not download or run
the models.

Within a cohort, identical inference content, or predictions within the
near-copy distance on the cohort batch, forms one cluster that receives one
credit split equally among its members (Section~\ref{reward-v1}). The
commitment nonce is excluded from content identity. The finalized commitment
block and extrinsic index, the admission order, break ties between equal
quantized scores and so choose which of the tied entrants heads a cluster and
can become the next reference. This order
assigns no authorship: it cannot establish who originally trained a model.
Already-published artifacts may enter but earn nothing
(Section~\ref{public-copies}). Permitted foundation-model components may be used
in new models, and a weight soup of published models, which fits the allowed
graph, is a new model.
